\subsection{Imaginary Roots Come in Pairs}
Do we know anything more about imaginary roots? First, let's recall a definition.
\begin{definition}[Complex Conjugate]
For any complex number $a+bi$, its \term{conjugate} is $a-bi$. (This is a symmetric relationship: the conjugate of $a+bi$ is $a-bi$, and the conjugate of $a-bi$ is $a+bi$.
\end{definition}
We know two things about complex conjugates.
\begin{itemize}
\item $(a+bi)(a-bi)=\makeblank[1in]$, so if we multiply a complex number by its conjugate we get a \makeblank.
\item $(a+bi)+(a-bi)=\makeblank[1in]$, so if we add a complex number to its conjugate we get a \makeblank.
\end{itemize}
This will be useful in understanding why the following theorem is true.
\begin{theorem}
For a polynomial with \emph{real} coefficients, if $a+bi$ is a root, its conjugate $a-bi$ is also.
\end{theorem}
This means that imaginary roots come in pairs. (If the root is real, then $b=0$, so $a+bi$ and $a-bi$ are the same number.) To see why, note that if we multiply
\[
\Bigl(x-(a+bi)\Bigr)\Bigl(x-(a-bi)\Bigr)
\]
\begin{lpic}{}{4}
\end{lpic}
The product has \makeblank[1in] coefficients.

(This doesn't prove the theorem---there might be some other way to get rid of the imaginary terms---but it does start to justify it.)